\subsection{Sparse DC fitting from aggregate geometry and a local row bound}
\label{restore:fit}

This subsection supplies the deterministic fitting module for the
degree-floor decoder. It uses the exact Bernoulli working criterion and
the exact-degree reconstruction $R^h$ in \eqref{restore:reconstruction}.
In particular, its rank-one degree correction and simple row hollowing
are different from the earlier correction $\diag(s)(H-I)$.
One profile refinement suffices. The optional finite repeated hard-refit
version is covered explicitly by \cref{restore:fit-invariant-basin}.

Put $d=np$, $B=pC$, $\Omega=\Theta ZBZ^\top\Theta$ and
$D_{0,i}=(\Omega\one)_i$. As before, $\mathbf d=A\one$ is the raw
degree vector, $D_i=\mathbf d_i$, and
\[
 G_i=\max(D_i,a_0\overline D),\qquad
 S=G^{-1/2}AG^{-1/2},\qquad
 R_0=G^{1/2}U\Lambda U^\top G^{1/2}.
\]
The constants $a_0>0$ and all subsequent clipping constants are fixed
using only the structural parameter bounds. They do not depend on $n$,
the unknown labels, or individual unobserved parameters.
Work on the two-sided total-degree event on which
$\overline D\asymp d$ and the population floor is inactive.
The aggregate geometric and block-total premises used below are
\begin{align}
 \norm{S-G^{-1/2}\Omega G^{-1/2}}_{\op}
   &\le C/\sqrt d,&
 \sum_iG_i&\le Cnd, \label{restore:fit-geometry-one}\\
 \sum_i\frac{(G_i-D_{0,i})^2}{G_i}
 +\sum_i\frac{(D_i-D_{0,i})^2}{G_i}
   &\le Cn,&
 \norm{(I-UU^\top)G^{1/2}Z}_F&\le C\sqrt n,
 \label{restore:fit-geometry-two}\\
 \max_{a,b}\left|[Z^\top(A-\Omega)Z]_{ab}\right|
   &\le Cn.&& \label{restore:fit-block-event}
\end{align}
Here $G_i\ge cd$, $D_{0,i}\asymp d$, and the retained signal has a fixed
gap from zero. The degree-floor proposition already supplies
\eqref{restore:fit-geometry-one}--\eqref{restore:fit-geometry-two}.
The fixed-$K$ Bernoulli block-total bound
\eqref{restore:fit-block-event} has failure probability
$C\exp\{-cL^2n/(d+L)\}$ when its fixed threshold constant is $L$.
Under $np^2=d^2/n\le C_1$, increasing $L$ pays this failure with any
prescribed fixed exponential rate in $d$. No degree-wise union bound
is needed for this block event.

Write $\Pi=UU^\top$. An aligned orthonormal basis of the random
population space has the exact form
\[
 V=G^{-1/2}\Theta Z
       (Z^\top\Theta G^{-1}\Theta Z)^{-1/2}O,
\]
where $O$ is orthogonal. The diagonal Gram matrix has entries comparable
to $n/d$. The signal gap and \eqref{restore:fit-geometry-one} imply
\begin{equation}\label{restore:fit-basis}
 \norm{U-V}_F\le C/\sqrt d,\qquad
 \sqrt{G_i}\norm{V_{i\cdot}}_2\le C\sqrt{d/n}.
\end{equation}
An orthogonal rotation used in this display does not change any row
norm or projection. Also $\norm S_{\op}\le1$: the inequality
$|x^\top Ax|\le\sum_iD_ix_i^2\le\sum_iG_ix_i^2$ proves this directly.
Hence $\norm\Lambda_{\op}\le1$.

\begin{lemma}[Uniform weighted mass of wrong-label sets]
\label{restore:fit-wrong-mass}
Fix $M<\infty$. For every possibly data-dependent set $E$ with
$|E|\le Mn/d$, on the aggregate event,
\begin{equation}\label{restore:fit-wrong-mass-bounds}
 \sum_{i\in E}G_i\le C_Mn,\qquad
 \sum_{i\in E}\sqrt{G_i}\norm{U_{i\cdot}}_2
       \le C_M\sqrt{n/d}.
\end{equation}
Moreover
$\sum_i\sqrt{G_i}\norm{U_{i\cdot}}_2\le C\sqrt{nd}$.
\end{lemma}

\begin{proof}
The elementary inequality
$G_i\le2D_{0,i}+(G_i-D_{0,i})^2/G_i$ gives
$\sum_{i\in E}G_i\le C(d|E|+n)\le C_Mn$.
Put $a_i=\sqrt{G_i}\norm{U_{i\cdot}}_2$.
By \eqref{restore:fit-basis} and Cauchy--Schwarz,
\[
 \sum_{i\in E}a_i
 \le C|E|\sqrt{d/n}
      +\left(\sum_{i\in E}G_i\right)^{1/2}\norm{U-V}_F
 \le C_M\sqrt{n/d}.
\]
The last assertion follows from
$\sum_i a_i\le(\sum_iG_i)^{1/2}\norm U_F$ and fixed $K$.
All estimates hold simultaneously for all such sets; they impose no
independence requirement on the selected labels.
\end{proof}

\begin{lemma}[Radius-count initialization under aggregate geometry]
\label{restore:fit-initializer}
The deterministic radius-count initializer specified earlier has at most
$C n/d$ errors on the aggregate event, after one common permutation.
This assertion requires neither a uniform degree upper bound nor
uniform spectral-row incoherence.
\end{lemma}

\begin{proof}
The normalized rows of $V$ are the $K$ orthogonal population centers.
Choose fixed $b_0>2\max_iD_{0,i}/d$, uniformly over the structural class.
For every $G_i>b_0d$,
$(G_i-D_{0,i})^2/G_i\ge G_i/4\ge b_0d/4$.
Thus at most $Cn/d$ such rows occur by
\eqref{restore:fit-geometry-two}. They contribute at most $4Cn/d$
to the squared directional error. For the remaining rows,
$\norm{V_{i\cdot}}_2\ge c/\sqrt n$. The row-normalization inequality
and \eqref{restore:fit-basis} bound their total squared directional
error by $Cn\norm{U-V}_F^2\le Cn/d$.
The existing radius-count lemma now applies with total squared error
$Cn/d$ and fixed positive community proportions. Its proof is
deterministic and unchanged.
\end{proof}

\begin{lemma}[Degree-corrected block totals and absolute mass]
\label{restore:fit-block-totals}
For $R^h$ as in \eqref{restore:reconstruction},
\begin{equation}\label{restore:fit-global-absolute}
 \sum_{i,j}|R^h_{ij}|\le Cnd.
\end{equation}
Uniformly over every fitted partition $\widehat Z$ having at most
$Mn/d$ errors under the common label matching,
\begin{equation}\label{restore:fit-total-accuracy}
 [\widehat Z^\top R^h\widehat Z]_{ab}
       =n_an_bB_{ab}+O_M(n).
\end{equation}
The omitted self-loops are included in the remainder.
\end{lemma}

\begin{proof}
Let $S_{\rm tail}=(I-\Pi)S(I-\Pi)$ and
$F_1=(I-\Pi)G^{1/2}\one$. Spectral invariance and the signal gap give
$\norm{S_{\rm tail}}_{\op}\le C/\sqrt d$ and
$\norm{F_1}_2\le C\sqrt n$.
Writing $v=\mathbf d-R_0\one$, one has exactly
$v=G^{1/2}S_{\rm tail}F_1$. Hence
\[
 \sum_i|v_i|
 \le\left(\sum_iG_i\right)^{1/2}
          \norm{S_{\rm tail}F_1}_2\le Cn.
\]
Since $\Pi_{ii}\le1$ and $\sum_i\Pi_{ii}=K$,
\[
 \sum_i|(R_0)_{ii}|
 \le\sum_iG_i\Pi_{ii}
 \le2(\max_iD_{0,i})K
       +\sum_i\frac{(G_i-D_{0,i})^2}{G_i}
 \le C(d+n)\le Cn.
\]
Consequently $\sum_i|(R_1)_{ii}|\le Cn$ and the simple hollowing
has total absolute cost at most $2\sum_i|(R_1)_{ii}|\le Cn$.
With $a_i$ as in the preceding lemma,
$|(R_0)_{ij}|\le a_ia_j$, so
$\sum_{i,j}|(R_0)_{ij}|\le(\sum_i a_i)^2\le Cnd$.
The rank-one correction has total absolute mass $\sum_i|v_i|\le Cn$.
This proves \eqref{restore:fit-global-absolute}.

Put $\widehat F=(I-\Pi)G^{1/2}\widehat Z$. By
\eqref{restore:fit-wrong-mass-bounds},
\[
 \norm{\widehat F}_F
 \le\norm{(I-\Pi)G^{1/2}Z}_F
       +\norm{G^{1/2}(\widehat Z-Z)}_F
 \le C_M\sqrt n .
\]
The exact quadratic identity
\[
 \widehat Z^\top(R_0-A)\widehat Z
       =-\widehat F^\top S_{\rm tail}\widehat F
\]
therefore has norm at most $C_Mn/\sqrt d$. Moving the mislabeled
vertices changes a raw block total by at most a fixed multiple of
$\sum_{i\in E}D_i\le\sum_{i\in E}G_i\le C_Mn$.
Equation \eqref{restore:fit-block-event} and the exact within-block
mean normalization of $\theta$ prove \eqref{restore:fit-total-accuracy}
with $R_0$ in place of $R^h$.
The rank-one correction changes a fitted block total by
$[\widehat Z^\top v]_a\widehat n_b/n$, where
$\widehat Z^\top v=\widehat F^\top S_{\rm tail}F_1=O_M(n/\sqrt d)$.
The already bounded hollowing cost completes the proof.
\end{proof}

\subsubsection{A bounded, normalized degree-parameter estimator}

Choose fixed positive $\alpha$ sufficiently small and $\beta$ sufficiently
large, using the structural bounds, so that on the two-sided
total-degree event
\[
 \alpha\overline D\le\tfrac12\min_iD_{0,i},\qquad
 \beta\overline D\ge2\max_iD_{0,i}.
\]
Define
\begin{equation}\label{restore:fit-clamped-theta}
 T_i=\min\{\max(D_i,\alpha\overline D),\beta\overline D\},
 \qquad
 \widehat\theta_i=
 \frac{T_i}{\widehat n_{\widehat z_i}^{-1}
                 \sum_{j:\widehat z_j=\widehat z_i}T_j}.
\end{equation}
This changes the nuisance estimator only. The retained degrees and the
degree identity $R^h\one=\mathbf d$ still use the original raw degrees.
The fitted parameters lie in $[\alpha/\beta,\beta/\alpha]$, and their
mean in every fitted block is exactly one.

\begin{lemma}[Clamped nuisance estimates]
\label{restore:fit-theta}
Uniformly over the $Mn/d$ count basin,
\begin{equation}\label{restore:fit-theta-rates}
 \frac{\norm{\widehat\theta-\theta}_2}{\sqrt n}
        \le C_Md^{-1/2},\qquad
 \widehat n_a^{-1}\sum_{\widehat z_i=a}T_i
        =dS_a+O_M(1),
\end{equation}
where $D_{0,i}=dS_{z_i}\theta_i$.
\end{lemma}

\begin{proof}
The chosen clamp interval contains $D_{0,i}$ with a fixed margin.
A direct split into the unclamped, lower-clamped, and upper-clamped
cases gives
\[
 |T_i-D_i|\le C\frac{(D_i-D_{0,i})^2}{G_i},\qquad
 |T_i-D_{0,i}|^2\le Cd\frac{(D_i-D_{0,i})^2}{G_i}.
\]
For the first inequality, an active lower clamp implies a degree
deviation of order $d$, and an active upper clamp implies a deviation
of order $D_i$. For the second, the unclamped range has $G_i\le Cd$;
projection onto an interval containing $D_{0,i}$ cannot increase
the deviation. Above the upper clamp,
$D_i((T_i-D_{0,i})/(D_i-D_{0,i}))^2\le T_i\le Cd$.
Summing yields
\begin{equation}\label{restore:fit-clamp-cost}
 \sum_i|T_i-D_i|\le Cn,\qquad
 \sum_i|T_i-D_{0,i}|^2\le Cnd.
\end{equation}
Since $\one=Z\one_K$, the raw block-total event implies
$\sum_{z_i=a}(D_i-D_{0,i})=O(n)$.
Clamping changes that sum by $O(n)$, and moving $Mn/d$ labels changes
a clamped degree total by at most $Cd(Mn/d)=O_M(n)$.
The fitted size differs from $n_a$ by $O_M(n/d)$.
Together with $\sum_{z_i=a}D_{0,i}=n_adS_a$, these facts prove the
group-mean assertion in \eqref{restore:fit-theta-rates}.
For correctly labeled vertices, division of
$T_i-D_{0,i}$ by this fitted mean, followed by
\eqref{restore:fit-clamp-cost}, gives squared error $O_M(n/d)$.
The incorrectly labeled vertices number $O_M(n/d)$ and have bounded
theta error. This proves the first assertion.
\end{proof}

\subsubsection{The exact Bernoulli block likelihood}

Fit a positive symmetric $B$ by maximizing the ordered Bernoulli
working criterion with $R^h$ and \eqref{restore:fit-clamped-theta},
divided by two, under the same fixed probability cap
$\widehat\theta_i\widehat\theta_jB_{\widehat z_i\widehat z_j}
\le q_{\max}<1$.
The next argument uses aggregate absolute mass instead of the stronger
and unnecessary assertion that every reconstructed entry is less than one.

\begin{lemma}[Unique Bernoulli block fit and its rate]
\label{restore:fit-block-root}
The fitted block roots are unique and interior for all large $n$.
Uniformly over the $Mn/d$ basin,
\begin{equation}\label{restore:fit-B-rate}
 \max_{a,b}\left|\frac{\widehat B_{ab}}{B_{ab}}-1\right|
       \le C_M(d^{-1}+p)\le C_Md^{-1}.
\end{equation}
\end{lemma}

\begin{proof}
For the ordered pairs $\mathcal J_{ab}$ associated with an unordered
block pair, put
\[
 E_{ab}=\sum_{\mathcal J_{ab}}R^h_{ij},\qquad
 t_{ij}=\widehat\theta_i\widehat\theta_j,\qquad
 T_{ab}=\sum_{\mathcal J_{ab}}t_{ij},\qquad
 \nu_{ab}=2-\mathbf1\{a=b\}.
\]
The preceding block-total lemma and exact fitted normalization imply
\[
 E_{ab}=\nu_{ab}n_an_bB_{ab}+O_M(n)\asymp n^2p>0,\qquad
 T_{ab}=\nu_{ab}\widehat n_a\widehat n_b+O(n).
\]
The exact score equation is $E_{ab}=F_{ab}(b)$, where
\[
 F_{ab}(b)=b\sum_{\mathcal J_{ab}}
              \frac{(1-R^h_{ij})t_{ij}}{1-bt_{ij}}.
\]
For every admissible $b$, bounded positive theta, the probability cap,
and \eqref{restore:fit-global-absolute} give
\[
 F_{ab}'(b)
 =\sum_{\mathcal J_{ab}}
          \frac{(1-R^h_{ij})t_{ij}}{(1-bt_{ij})^2}
 \ge cn^2-C\sum_{\mathcal J_{ab}}|R^h_{ij}|
 \ge cn^2-Cnd>0.
\]
Here $p=d/n\to0$. At zero, $F_{ab}(0)=0$; its value at the
positive cap endpoint is of order $n^2$, while $E_{ab}\asymp n^2p$.
There is a unique interior root. The likelihood derivative is
$(E_{ab}-F_{ab}(b))/b$, so the root is its unique global maximum.

For $b\asymp p$,
\[
 \sum_{\mathcal J_{ab}}
       \frac{(1-R^h_{ij})t_{ij}}{1-bt_{ij}}
       =T_{ab}+O(n^2p).
\]
The error follows by expanding the denominator and using
$\sum|R^h_{ij}|\le Cnd=Cn^2p$.
Bracketing the root between fixed multiples of $p$, then using this
expansion, gives
$\widehat B_{ab}=(E_{ab}/T_{ab})\{1+O(p)\}$.
The block-total and group-size errors prove the first inequality in
\eqref{restore:fit-B-rate}. The second uses $dp=np^2\le C_1$.
\end{proof}

\subsubsection{A local row condition suffices for every fitted stage}

For a tested vertex $i$, impose only the local conditions
\begin{equation}\label{restore:fit-local-row}
 D_i\le bd,\qquad
 a_i:=\sqrt{G_i}\norm{U_{i\cdot}}_2\le L\sqrt{d/n},
\end{equation}
where $b,L$ are fixed. These are not imposed on every neighbor
simultaneously.

\begin{lemma}[Local absolute mass and wrong-label stability]
\label{restore:fit-local-mass}
On the aggregate event and \eqref{restore:fit-local-row},
\begin{equation}\label{restore:fit-local-mass-bounds}
 \sum_j|R^h_{ij}|\le C_{b,L}d,\qquad
 \sup_{|E|\le Mn/d}\sum_{j\in E}|R^h_{ij}|\le C_{M,b,L}.
\end{equation}
\end{lemma}

\begin{proof}
Since $|(R_0)_{ij}|\le a_ia_j$,
\[
 \sum_j|(R_0)_{ij}|\le a_i\sum_j a_j\le Cd,\qquad
 \sum_{j\in E}|(R_0)_{ij}|
       \le a_i\sum_{j\in E}a_j\le C_M
\]
by \cref{restore:fit-wrong-mass}. Also
$|v_i|\le D_i+\sum_j|(R_0)_{ij}|\le Cd$.
The rank-one correction therefore adds at most $Cd$ to the first
bound and $|v_i||E|/n\le C_M$ to the second.
Finally $|(R_1)_{ii}|\le a_i^2+|v_i|/n\le Cp$.
The simple row hollowing changes either row bound by at most
$2|(R_1)_{ii}|\le Cp$. This proves the claim.
\end{proof}

\begin{lemma}[Bernoulli profile reduction for a signed row]
\label{restore:fit-profile}
Let $r$ be a hollow signed row satisfying
$\sum_jr_j=D_i\ge0$ and $\sum_j|r_j|\le Ld$.
Let the positive neighbor intensities $b_j$ lie between fixed
multiples of $p$. Maximizing the exact Bernoulli row criterion
\[
 \sum_j\{r_j\log(tb_j)+(1-r_j)\log(1-tb_j)\}
\]
over $t\ge0$ under the fixed cap $tb_j\le q_{\max}<1$ gives the
usual normalized block-profile contrast, with remainder $O(np^2)$.
No entrywise bound on $r_j$ is required.
\end{lemma}

\begin{proof}
After combining the coefficient of $\log t$, the derivative is
\[
 D_i/t-\sum_j\frac{(1-r_j)b_j}{1-tb_j}.
\]
Throughout the admissible interval,
\[
 \sum_j\frac{(1-r_j)b_j}{1-tb_j}
 \ge cnp-Cp\sum_j|r_j|
 \ge cd-CLdp\ge c'd.
\]
Thus every maximizing $t$ lies in a fixed bounded interval,
because $D_i\le Ld$. This reasoning does not require individual
nonedge weights $1-r_j$ to be nonnegative.
On that fixed interval, the exact Bernoulli expansion is uniform:
\[
 \ell_r(t)
   =D_i\log t+\sum_jr_j\log b_j
         -t\sum_jb_j+O(np^2).
\]
Restoring the edge terms in the linear exposure costs
$Cp\sum_j|r_j|=O(np^2)$, and the quadratic remainder costs
$Cp^2(n+\sum_j|r_j|)=O(np^2)$.
The leading optimizer $D_i/\sum_jb_j$ is bounded and admissible
for all large $n$, so taking maxima preserves the remainder.
The resulting label-dependent maximum is the standard normalized
block-profile score. If $D_i=0$, the derivative bound makes $t=0$
optimal after combining the $\log t$ terms; use the continuous
convention $0\log0=0$.
\end{proof}

\begin{proposition}[Uniform replacement by the true DC contrast]
\label{restore:fit-uniform-score}
\label{restore:fit-score}
Fix $M,b,L$. Use the clamped theta estimates
\eqref{restore:fit-clamped-theta}, the exact Bernoulli block fit, and
the exact Bernoulli target profile on $R^h$.
On the aggregate event and the local target condition
\eqref{restore:fit-local-row}, uniformly over every fitted partition
in the $Mn/d$ basin,
\begin{equation}\label{restore:fit-true-score}
 \widehat L_i^{ab}
       =[R^h w_{i,ab}]_i+O_{M,b,L}(1),\qquad
 (w_{i,ab})_j=
     \log\frac{q_{a,z_j}^{(i)}}{q_{b,z_j}^{(i)}}\quad(j\ne i).
\end{equation}
The true profiles $q^{(i)}$ are the self-excluded neighbor profiles
defined earlier. The value of $w_{i,ab}$ at $i$ is immaterial since
$R^h_{ii}=0$.
\end{proposition}

\begin{proof}
By \cref{restore:fit-profile}, the profile reduction costs
$O(np^2)=O(1)$. Exact fitted-block normalization gives the fitted
neighbor mass
$\widehat M_c^{(i)}
 =\widehat n_c-\widehat\theta_i\mathbf1\{\widehat z_i=c\}$.
The fitted sizes differ from the true sizes by $O_M(n/d)$,
the fitted theta values are bounded, and
\eqref{restore:fit-B-rate} gives relative block error $O_M(1/d)$.
Smoothness of normalization and logarithms on the fixed positive
profile class therefore yields fitted log-profile error $O_M(1/d)$.
Multiplication by the row absolute mass in
\eqref{restore:fit-local-mass-bounds} costs $O_M(1)$.
Replacing the fitted labels by the true labels changes a bounded
coefficient vector on at most $Mn/d$ coordinates; the second bound
in \eqref{restore:fit-local-mass-bounds} gives another $O_M(1)$.
All inequalities hold simultaneously for every partition in the basin,
so data-dependent fitting does not introduce an independence assumption.
\end{proof}

\begin{lemma}[A uniform basin for finite hard refits]
\label{restore:fit-invariant-basin}
\label{restore:fit-invariant}
There is a fixed $M$ such that, on the aggregate event and for all
sufficiently large $n$, each hard nuisance-fit and profile-assignment
update maps the $Mn/d$ error basin into itself under the same label
matching. Outside a fixed target-exception set determined by the graph,
every error at every updated iterate is contained in
\[
 [R^h w_{i,ab}]_i\le C_M
 \quad\text{for some }b\ne z_i.
\]
The target-exception set can be taken to be the failure of the local
conditions \eqref{restore:fit-local-row}, with fixed sufficiently large
thresholds. Thus a finite number of these specified hard refits, including
adaptive finite stopping after at least one update, requires no union
over iterations in the final risk bound.
\end{lemma}

\begin{proof}
Fix sufficiently large constants $b,R$. The set
$\{i:D_i>bd\}$ has at most $Cn/d$ elements by the degree energy in
\eqref{restore:fit-geometry-two}, since $b$ exceeds twice the largest
population degree multiplier. Since $G_i\ge cd$,
\eqref{restore:fit-basis} gives $\norm{V_{i\cdot}}_2\le C/\sqrt n$.
For $R>2C$, the set $\{i:\sqrt n\norm{U_{i\cdot}}_2>R\}$
also has at most $Cn/d$ elements, by
$\norm{U-V}_F^2\le C/d$.
Every other target satisfies \eqref{restore:fit-local-row} for a
fixed $L$, because $G_i\le C_bd$ there.
This target-exception set depends only on the graph and its
cardinality bound is independent of $M$.

Write $M_G=G^{-1/2}\Omega G^{-1/2}$ and
$S_K=U\Lambda U^\top$. Best rank-$K$ approximation and
\eqref{restore:fit-geometry-one} imply
$\norm{S_K-M_G}_{\op}\le C/\sqrt d$. Hence
\[
 \norm{G^{-1/2}(R_0-\Omega)Z}_F^2
 \le \frac Cd\norm{G^{1/2}Z}_F^2\le Cn.
\]
Moreover
$\norm{G^{-1/2}v}_2^2
 =\norm{S_{\rm tail}F_1}_2^2\le Cn/d$.
The rank-one correction therefore preserves the bound
\begin{equation}\label{restore:fit-basin-action}
 \norm{G^{-1/2}(R_1-\Omega)Z}_F^2\le Cn.
\end{equation}
The simple hollowing changes any bounded block-constant contrast at
row $i$ by at most $C|(R_1)_{ii}|$.
Since $\sum_i|(R_1)_{ii}|\le Cn$, at most $Cn/d$ rows can incur
a hollowing contrast error larger than any prescribed fixed small
multiple of $d$.
This auxiliary counting set is needed only for the deterministic
basin proof; it is not added as a probability-cost event in the sharp
risk argument.

The true profile contrast has a uniform positive population margin
$[\Omega w_{i,ab}]_i\ge c_{\rm gap}d$ for $a=z_i$, $b\ne a$,
including its $O(p)$ self-exclusion correction.
Extend each $w_{i,ab}$ to a block-constant vector at coordinate $i$;
this does not affect $R^h w$ because $R^h_{ii}=0$.
Its $K$ block coefficients are uniformly bounded. For a regular
target, \cref{restore:fit-uniform-score} bounds the fitted discrepancy
by $C_M$, which is less than $c_{\rm gap}d/4$ for all sufficiently
large $n$ (the threshold may depend on $M$).
Outside the auxiliary hollowing set, a new erroneous label therefore
forces
\[
 \norm{[(R_1-\Omega)Z]_{i\cdot}}_2\ge c d.
\]
For the regular targets, $G_i\le C_bd$. Summing the squared
requirement divided by $G_i$ and applying
\eqref{restore:fit-basin-action} gives at most $C_0n/d$ such targets.
Including the two cardinality exceptions still gives $C_0n/d$
with $C_0$ independent of $M$.

Choose $M$ at least both this $C_0$ and the initializer constant.
The basin is invariant for all sufficiently large $n$.
For every regular target, the canonical error-event inclusion
follows directly from \eqref{restore:fit-true-score}, independently
of this counting argument. It is the same event for all partitions
in the fixed basin. Induction therefore covers each finite hard
refit and any data-dependent finite stopping rule without a new
probability union or a stopping-time independence claim.
\end{proof}

\subsubsection{Ordinary SBM with a block Bernoulli fit and no target profiling}

The sparse ordinary-SBM decoder in this subsection sets $\theta_i=1$,
uses the same radius-count initializer and the same $R^h$, fits a
symmetric block Bernoulli matrix, and compares the full Bernoulli row
likelihoods. The target parameter is not profiled. This is an explicitly
specified block-fitting decoder; it is not a claim about the earlier
unmodified free-profile mixture E/M schedule.

\begin{proposition}[Ordinary block fitting and the true Bernoulli LRT]
\label{restore:fit-ordinary}
Assume ordinary SBM with a fixed positive well-conditioned block profile,
comparable group sizes, $d=np\to\infty$ and $np^2\le C_1$.
On the aggregate event, the radius-count initialization has $O(n/d)$
errors. For every fitted partition in a fixed $Mn/d$ basin, the symmetric
block Bernoulli working MLE based on $R^h$ satisfies
\begin{equation}\label{restore:fit-ordinary-parameter}
 \max_{a,b}|\widehat P_{ab}-P_{ab}|\le C_M/n.
\end{equation}
At every target obeying \eqref{restore:fit-local-row}, its full
Bernoulli contrast satisfies
\begin{equation}\label{restore:fit-ordinary-score}
 \widehat L_i^{ab}
    =[R^h w_{ab}]_i+c_i^{ab}+O_M(1),
\end{equation}
where
\[
 (w_{ab})_j=
   \log\frac{P_{a,z_j}(1-P_{b,z_j})}
                  {P_{b,z_j}(1-P_{a,z_j})},\qquad
 c_i^{ab}=\sum_{j\ne i}
       \log\frac{1-P_{a,z_j}}{1-P_{b,z_j}}.
\]
The counterpart of \cref{restore:fit-invariant-basin} also holds for
these ordinary block fits and full Bernoulli assignments.
\end{proposition}

\begin{proof}
For each unordered fitted block pair, let $E_{ab}$ be the corresponding
ordered reconstructed total and $N_{ab}$ its number of off-diagonal
ordered pairs. The exact Bernoulli working criterion is
$E_{ab}\log b+(N_{ab}-E_{ab})\log(1-b)$.
The block-total lemma gives
$E_{ab}=\nu_{ab}n_an_bP_{ab}+O_M(n)$ and
$N_{ab}=\nu_{ab}\widehat n_a\widehat n_b+O(n)$.
For all sufficiently large $n$, $0<E_{ab}<N_{ab}$; its unique MLE is
$\widehat P_{ab}=E_{ab}/N_{ab}$.
Any fixed sufficiently small probability floor of order $\overline D/n$
is inactive on this event. Since
$|\widehat n_a-n_a|\le Mn/d$ and $P_{ab}\asymp p$,
\[
 |\widehat P_{ab}-P_{ab}|
 \le C_M\{n/n^2+p(n/d)/n\}=C_M/n.
\]
This proves \eqref{restore:fit-ordinary-parameter}.

Writing the Bernoulli contrast as an incident-edge log-odds action
plus its nonedge intercept is exact. On the fixed positive profile
class, \eqref{restore:fit-ordinary-parameter} changes each log-odds
coefficient by $O_M((np)^{-1})=O_M(d^{-1})$.
Its total contribution at a regular target is $O_M(1)$ by the row
absolute-mass bound. Incorrect fitted neighbor labels change a bounded
log-odds coefficient on at most $Mn/d$ coordinates, at a total row
cost $O_M(1)$ by \cref{restore:fit-local-mass}.

For the intercept, the derivatives of $\log(1-b)$ are bounded on
the relevant probability range. Parameter replacement therefore costs
$n\,O_M(1/n)=O_M(1)$. Changing $Mn/d$ group assignments changes
each nonedge log-ratio by at most $Cp$, so its additional cost is
$C_M(n/d)p=O_M(1)$. These observations prove
\eqref{restore:fit-ordinary-score} for all data-dependent basin
partitions at once.

For the basin argument, the true population contrast equals a sum
of Bernoulli Kullback--Leibler divergences, up to the $O(p)$ completed
diagonal, and is at least $c_{\rm gap}d$.
The block-constant log-odds vectors are uniformly bounded.
Consequently the weighted action estimate
\eqref{restore:fit-basin-action}, the same target-exception count,
and the same hollowing count prove an output error bound $C_0n/d$.
As before $C_0$ is independent of the chosen input basin radius $M$;
only the sufficiently-large-$n$ threshold depends on $M$.
Choosing $M$ large enough proves invariance. Every regular-target
error at every refit is contained in the same true-contrast event
\eqref{restore:fit-ordinary-score} with its fixed $O_M(1)$ threshold,
so finite adaptive hard refits need no iteration union.
\end{proof}

The remaining interface is local. The radius-count initializer and
nuisance fit above need only aggregate geometry. To turn
\eqref{restore:fit-true-score} into the actual sharp risk theorem one
pays the target-local spectral-row exception, then applies the local
full-to-LOO comparison in \cref{restore:loo-corrected} to the deterministic
true contrast. The fitted coefficients themselves are not conditioned
on as if independent of the tested row.
The high-degree part of \eqref{restore:fit-local-row} has arbitrarily
large fixed exponential rate after increasing $b$. On $D_i\le bd$,
the spectral-row part is controlled by the truncated-row estimate,
since $G_i\le C_bd$ and
$a_i>L\sqrt{d/n}$ implies
$\sqrt n\norm{U_{i\cdot}}_2>L/\sqrt{C_b}$.
